% =====================================================
% Method — target length ~1.5-2 pages, per SPEC.md's
% "Paper writing plan (Step 6.2) -- narrative structure".
%
% Every number, formula and design fact below is taken from
% SPEC.md, PREREGISTRATION.md, src/qgemm/bounds.py,
% src/qgemm/metrics.py and configs/sweep_main.yaml. No results
% are previewed here (Steps 4.1-4.4 belong to Section 4).
%
% Notation note: SPEC.md calls the derived bound `cota(n)`.
% Rendered here as \beta(n) -- an English-language paper should
% not carry the Spanish word untranslated. Same object, renamed
% only for the write-up.
%
% Base LaTeX math only (\[ \], equation, \frac, \sqrt, \mathrm):
% main.tex does not load amsmath, and this file deliberately does
% not require it.
% =====================================================

\section{Method}
\label{sec:method}

\subsection{Setup: backward error under exact accumulation}
\label{sec:method:setup}

The measurement instrument is a quantized GEMM. For operands $A \in
\mathbb{R}^{M \times n}$ and $B \in \mathbb{R}^{n \times N}$, both are quantized
with a block-scaled low-precision format along the contraction dimension $n$,
and the resulting float64 \emph{reconstructions} --- the dequantized values, not
the integer codes --- are multiplied. Accuracy is reported elementwise as the
backward error \citep{rasquinha2023metric}
\begin{equation}
  \mathrm{BE}_{ij} \;=\; \frac{\bigl|(AB)_{ij} - \hat{C}_{ij}\bigr|}
                              {\bigl(|A|\,|B|\bigr)_{ij}},
  \label{eq:be}
\end{equation}
with $\hat{C}$ the quantized product and $|\cdot|$ elementwise. The denominator
uses the product of absolute values, not $|AB|$; the two differ whenever the
exact product involves cancellation.

The primary route accumulates the reconstructions \emph{exactly}, in float64 ---
a design choice, not an implementation convenience. A low-precision GEMM has two
independent error sources, operands snapped onto a coarse grid and partial sums
rounded in a narrow accumulator, and they compose, so a number measured with
both active cannot be attributed to either. Under exact accumulation the only
remaining approximation is quantization of the inputs, so the residual
$\hat{C} - AB$ \emph{is} the input-quantization error and (\ref{eq:be}) isolates
it. A bfloat16 accumulator exists as a deliberate secondary ablation, outside
the grid analysed here and never mixed into a headline result.

Operands have i.i.d.\ t-Student($\nu$) entries --- smaller $\nu$ meaning heavier
tails, $\nu \to \infty$ the Gaussian case --- and each tensor is normalized as a
whole to unit median absolute deviation, rather than unit standard deviation,
which is undefined for the $\nu \le 2$ part of the grid. Every cell uses the
shape $(64,n) \times (n,64)$.

\subsection{Why the classical bound does not transfer}
\label{sec:method:classical}

The classical worst-case backward-error bound for a length-$n$ inner product is
$\gamma_n = nu/(1-nu)$ at unit roundoff $u$ \citep{higham2002accuracy};
probabilistic analyses relax the leading factor to $\sqrt{n}\,u$, or
$\sqrt{n \log n}\,u$ for a refined variant, as a high-probability upper bound
rather than a claim about typical growth
\citep{higham2019new,connolly2021stochastic}.

Neither transfers here, for a mechanistic reason. The $n$-dependence of both
families arises specifically from \emph{repeated rounding during accumulation}:
one rounding per summation step, $n$ of them, compounding either in the worst
case ($\gamma_n$) or as a random walk ($\sqrt{n}\,u$). Under exact accumulation
that mechanism is not physically present in the pipeline, so whatever
$n$-dependence the measured error has, it cannot be the one those bounds
describe. A block-scaled format also offers no fixed $u$ to substitute: the
classical unit roundoff is half an ulp of a grid fixed by the encoding, whereas
the grid an element lands on here is set by the largest magnitude in \emph{its
own block}, making the achieved relative error a distribution rather than a
property of the format.

\subsection{A measured unit roundoff, and the bound built from it}
\label{sec:method:bound}

The quantity replacing $u$ is measured directly. For a given element format,
block size, scale format and tail index $\nu$, define
\begin{equation}
  u_{\mathrm{eff}}(\mathrm{format}, \mathrm{block}, \nu)
  \;=\; \mathrm{median}_{\,i}\;
        \frac{\bigl|x_i - \hat{x}_i\bigr|}{\bigl|x_i\bigr|},
  \label{eq:ueff}
\end{equation}
the median per-element relative quantization error: draw samples from that
$\nu$, quantize under that configuration, read the median off the resulting
distribution. Its $99$th percentile is recorded alongside. Estimation uses an
independent calibration draw, seeded from a separate stream, so the denominator
of the ratio tested below is never calibrated on the sample it is tested against.

One choice inside (\ref{eq:ueff}) is load-bearing. Relative error is ill-defined
as $x_i \to 0$, and ``near zero'' means something only relative to an element's
own effective scale $s_{\mathrm{eff}}$ (block scale times global scale, where
one is used), since block scales span many orders of magnitude across a
heavy-tailed tensor. Elements with $|x_i| \le 0.25\,s_{\mathrm{eff}}$ are
therefore excluded. The threshold is not tuned: E2M1's smallest nonzero
magnitude is $0.5\,s_{\mathrm{eff}}$, so under round-to-nearest every element at
or below $0.25\,s_{\mathrm{eff}}$ quantizes to exactly zero with relative error
exactly $1$, recording the format's dynamic-range floor rather than its
precision. That the cut is necessary was checked, not assumed: with it disabled,
all $32$ measured configurations return $p_{99} = 1.00000$ exactly. Its price is
a large, configuration-dependent exclusion --- surviving fractions range from
$0.360$ (MXFP4 at $\nu=1$) to $0.932$ (NVFP4 at the Gaussian end) --- tabulated
wherever $u_{\mathrm{eff}}$ is reported.

Under exact accumulation the error at output entry $(i,j)$ is a sum of $n$
one-time input-quantization errors of essentially random sign, so partial
cancellation makes it grow like $\sqrt{n}$ rather than $n$, while the
denominator of (\ref{eq:be}), $\sum_k |A_{ik}|\,|B_{kj}|$, sums $n$ strictly
positive terms and grows approximately linearly for well-behaved inputs.
Together these give $\mathrm{BE} \sim \sqrt{n}\,u_{\mathrm{eff}}/n$, and hence
the tested bound
\begin{equation}
  \beta(n, \mathrm{format}, \mathrm{block}, \nu)
  \;=\; c \cdot
  \frac{u_{\mathrm{eff}}(\mathrm{format}, \mathrm{block}, \nu)}{\sqrt{n}},
  \label{eq:bound}
\end{equation}
which \emph{decays} in $n$ instead of growing. This is an order-of-magnitude
argument of the central-limit\slash law-of-large-numbers kind, not a proven
tight inequality, and is not offered as a theoretical contribution: both
premises --- a mean-zero numerator averaging like a random walk, a denominator
obeying a law of large numbers --- require moments the heaviest tails swept here
lack, and at $\nu=1$ (Cauchy) the derivation is expected to break down. That is
a limit of its assumptions, not a defect in the form of (\ref{eq:bound}).

The constant $c$ is fixed a priori at $c=1$ for all confirmatory analysis, per
the anti-circularity principle of this study's preregistration: a bound whose
constant is fitted to the same data it is tested against cannot be falsified by
that data. Keeping a unit leading constant follows the precedent of
probabilistic rounding analyses, which substitute $\sqrt{n}\,u$ for the
deterministic factor $nu$ without introducing a fitted constant
\citep{connolly2021stochastic}. Since the derivation is order-of-magnitude
rather than tight, a poorly calibrated $c=1$ is itself a reportable finding
rather than grounds to re-fit; any $\hat{c}$ estimated from data stays
descriptive and exploratory, never substituted back into the confirmatory ratio.

The bound is declared broken at $(\mathrm{format},\nu,n)$ when the lower end of
the $95\%$ bootstrap confidence interval of the median of $\mathrm{BE}/\beta(n)$
exceeds $1.0$: when error decays reliably more slowly than the cancellation rate
predicts, not when it exceeds a growing worst-case ceiling. Because
$u_{\mathrm{eff}}$ is itself indexed by $\nu$, both sides of the ratio move with
$\nu$, so this tests whether the \emph{scaling law} survives heavy tails, not
whether heavy tails increase absolute error.

\subsection{Factorial design}
\label{sec:method:design}

MXFP4 and NVFP4 differ in \emph{both} block size and scale format at once, so
any difference measured between the presets is unattributable --- block, scale
format, or interaction. The design crosses the two factors, block size $\in
\{16,32\}$ against scale format $\in \{\mathrm{E8M0},\mathrm{E4M3}\}$, so either
can be held fixed while the other varies; this $2 \times 2$ is the causal core.
The element format is E2M1 everywhere, since both real formats store E2M1 and
fixing it makes the scaling the only thing that varies. Use of a per-tensor
global scale is derived from the scale format rather than swept: E8M0 spans
$255$ binades and reaches any block's magnitude unaided, while E4M3 spans about
$19$ and overflows to NaN without one.

Two of the four combinations are real formats: block $32$ with an E8M0 scale is
MXFP4, block $16$ with an E4M3 scale is NVFP4. \textbf{The other two --- block
$16$ with E8M0, and block $32$ with E4M3 --- are experimental controls.} They
correspond to no hardware, no specification and no vendor format, and exist
solely as measurement instruments that make the comparison between the two real
formats separable. Nothing here proposes, endorses or benchmarks them as
formats, and results tables keep the distinction visible.

Rounding mode (round-to-nearest-even or stochastic) and a random Hadamard
transform (on or off) are secondary factors, not part of the causal core. The
transform is applied per block along the contraction dimension with the same
random signs on both operands, so it cancels inside every inner product and
needs no inversion, spreading a block's energy so a lone outlier no longer sets
the block scale. Per block by design: a global, whole-row transform spreads by
$\sqrt{n}$, tying its effect to the same $n$ whose scaling law this study sweeps.

These four factors are crossed with the tail index, which runs over
$\nu \in \{1,2,3,5,8,15,30\}$ together with the Gaussian limit, and with the
contraction dimension $n \in \{16,64,256,1024,4096\}$, for a grid of
$8 \times 5 \times 2^4 = 640$ cells, all with exact accumulation. The primary contraction
dimension is fixed in advance at $n=4096$, closest to real GEMM shapes in
language-model workloads; the other four are robustness checks, reported as a
disclosed sensitivity table and never substituted for it. Trial counts are
adaptive: $5000$ per cell at $\nu \in \{1,2,3\}$, where quantile estimates are
noisiest and the bound is most likely to break, and $1000$ per cell for the
lighter tails, for a total of $1.6$ million trials. Three per-tensor formats
(FP8-E4M3, FP8-E5M2 and INT8) are swept over the same $\nu$ grid at
$n \in \{16,256,4096\}$ as context baselines, outside the factorial and outside
every confirmatory comparison.

\subsection{Robust statistical protocol}
\label{sec:method:stats}

Each trial contributes a summary of its own $M \times N$ backward-error matrix,
and the \emph{trial} is the resampling unit throughout --- never the individual
matrix element, since elements within a block share a scale factor, are not
exchangeable, and would give anticonservative intervals that inflate the
measured break rate. Confidence intervals use the percentile bootstrap,
$B = 10{,}000$ resamples, an explicitly seeded generator, and are reported as
computed rather than forced symmetric.

The median and $90$th percentile carry all confirmatory statistical weight; the
$99$th percentile is reported for every cell but is indicative only, at any
trial count used here. The justification is a measured coverage gap, not a
generic warning that percentile bootstraps can be unreliable. A cell's $p_{99}$
is set by roughly its top $1\%$ of values ($10$ of $1000$ trials, $50$ of
$5000$), and a percentile bootstrap resamples with replacement from the observed
sample, so it can reweight those values but never produce one more extreme than
the largest already drawn. A coverage simulation quantifies this --- $1000$
simulated experiments per cell, t-Student($\nu=2$), both trial budgets, nominal
$95\%$: median coverage is $95.5\%$ at $1000$ trials and $94.8\%$ at $5000$,
close to nominal at both, while $p_{99}$ coverage is $92.7\%$ and $94.0\%$,
below the median's at both budgets and roughly $3.5$ standard errors below
nominal at $1000$ trials. The gap is real and directionally as the mechanism
predicts, but modest rather than a collapse; more trials measurably help
$p_{99}$ without closing it, which is why the indicative-only status holds at
every trial count rather than below some threshold.
